% Cuerpo íntegro, sin preámbulo ni portada duplicada.



\section*{Objeto y resultados}
Se demuestra que el refinamiento reversible con acarreo producido por la
arquitectura APP--TRIT--TPK admite una forma canónica estable: la suma de las
formas de las celdas refinadas coincide exactamente con la forma de la celda
madre y todos los polos interiores se cancelan. La identidad se extiende de
manera exacta al complejo producto cuboidal y conserva los residuos iterados de
la frontera externa. El registro dodecafásico \(K\) determina un marco
invertible \(O_K\), una métrica de Gram y un archivo isométrico que preservan
volúmenes exteriores y permiten recuperar la historia. Un cálculo exhaustivo
de menores muestra que \(O_K\) no es totalmente no negativo. La dirección
ya construida \(u_K=P_3P_{11}K\), combinada con
\(\rho_A=\pi_{\rm HMT}/\varphi_{\rm HMT}^2\), define un flujo diagonal
positivo de determinante uno que conserva cada signo de Plücker y cada
estrato positroidal. Esta separación permite formular, con
condiciones falsables, el atlas positroidal intrínseco del TPK, su dinámica de
cúmulos y el empuje hacia una sección amplituhedral. La centralidad del
electrón queda vinculada al único punto fijo \((5,5)\) del atlas APP, sin
confundirlo con el dato nonádico \(111111\) ni con la firma regional
\(555555\).


\section{Tesis y orden causal}

El objeto de partida no es un Grassmanniano elegido exteriormente. La cadena
\[
\APP\longrightarrow\mathrm{TRIT}\longrightarrow\TPK
\longrightarrow\widehat x
\longrightarrow\mathcal C_{\rm cont}^{\rm disc}
\]
produce un estado enriquecido que conserva hojas, orientación, residuo,
cociente, acarreo, ruta, frontera, supervivencia, memoria y retorno. Las cinco
construcciones del continuo emergen conjuntamente. La geometría positiva es
una publicación de su refinamiento compatible.

El artículo separa tres afirmaciones que no deben fusionarse:
\begin{enumerate}
\item la telescopía y el archivo métrico son resultados exactos;
\item la curva de momentos, \(D_K(s)\) y el flujo centrado
      \(g_{K,\rho}(\tau)\) son formalizaciones nuevas exactas;
\item el atlas positroidal intrínseco, la semilla de cúmulo y la igualdad con
un amplituhedro completo exigen mapas adicionales expresamente declarados.
\end{enumerate}
Esta separación no disminuye el resultado: identifica la prueba de cada flecha
y evita atribuir a una matriz métrica una positividad que no posee.

\section{Teorema matriz}


% BEGIN OWNER /Users/ruben/Documents/New project/output/IMPLEMENTACION_GEOMETRIA_POSITIVA_TPK_20260918/02_TEOREMA_TRONCAL_GEOMETRIA_POSITIVA_TPK.tex
% Cuerpo de inserción autónomo. No modifica ningún manuscrito activo.
\section{Geometría positiva coinductiva del refinamiento TPK}
\label{geometria_positiva:sec:geom-positiva-tpk}

La estructura positiva que se estudia en esta sección no se introduce desde
un Grassmanniano previamente supuesto. Procede de la misma cadena generadora
que construye el estado enriquecido:
\[
 \APP\longrightarrow\mathrm{TRIT}\longrightarrow\TPK
 \longrightarrow \widehat x
 \longrightarrow\mathcal C_{\rm cont}^{\rm disc}.
\]
APP conserva simultáneamente sus dos hojas, sus residuos y sus cocientes;
TRIT tipa el régimen y orienta cada transición; TPK selecciona, transporta y
actualiza sin perder acarreo, ruta, frontera, supervivencia ni memoria. Las
cinco construcciones del continuo emergen conjuntamente de ese estado. La
realización grassmanniana que sigue es un lector posterior de su refinamiento
compatible y no una entrada que seleccione la dinámica.

\subsection{Acarreo, subdivisión y parámetro ordenado}

Fijemos una base entera \(B\ge2\). La operación elemental de refinamiento con
acarreo \(c\in\{0,\ldots,B-1\}\) es
\begin{equation}
 L_c(x,y)=(Bx-c,By+c).
 \label{geometria_positiva:eq:carry-map-positive}
\end{equation}
Conserva la suma reescalada, pues
\[
 \frac{Bx-c}{B(x+y)}+\frac{By+c}{B(x+y)}=1.
\]
Su inversa recupera \(c\) por la clase residual declarada y después recupera
\((x,y)\) mediante división euclídea. Por composición, una palabra de acarreos
\(c_1\ldots c_n\) determina una celda de profundidad \(n\), y las divisiones
en orden inverso recuperan la palabra completa. La compatibilidad entre
truncamientos produce cilindros encajados y conserva la historia, no sólo su
evaluación límite.

Sea \(M\ge1\) el número de intervalos de la carta inicial y
\begin{equation}
 N_n=B^nM,\qquad t_{n,j}=\frac{j}{N_n},\qquad 0\le j\le N_n.
 \label{geometria_positiva:eq:ordered-carry-points}
\end{equation}
Para \(0\le j<N_n\), el subintervalo de índice \(c\) del punto \(j\) tiene índice
\begin{equation}
 j'=Bj+c,\qquad
 t_{n+1,Bj+c}=t_{n,j}+\frac{c}{BN_n}.
 \label{geometria_positiva:eq:carry-insertion}
\end{equation}
La desigualdad \(j<j'\Rightarrow t_{n,j}<t_{n,j'}\) hace visible la
información que necesita la positividad: el acarreo no sólo refina; inserta
puntos conservando el orden total compatible con los mapas de enlace.

\subsection{Matriz de momentos y Grassmanniano positivo}

\paragraph{Estatuto.}
La curva de momentos que sigue es una formalización nueva del orden ya
producido por el refinamiento; no se presenta como un operador adicional del
TPK ni como una entrada del generador.

Para \(1\le r\le N_n+1\), definimos la curva de momentos y su matriz de
columnas
\[
 v_r(t)=(1,t,\ldots,t^{r-1})^{\mathsf T},\qquad
 Z_n^{(r)}=(v_r(t_{n,0})\ \cdots\ v_r(t_{n,N_n})).
\]

\begin{theorem}[Positividad de Plücker generada por el acarreo]
\label{geometria_positiva:thm:tpk-vandermonde-positive}
Para todo \(I=\{j_1<\cdots<j_r\}\),
\begin{equation}
 \Delta_I(Z_n^{(r)})
 =\prod_{1\le a<b\le r}(t_{n,j_b}-t_{n,j_a})>0.
 \label{geometria_positiva:eq:vandermonde-minor}
\end{equation}
Por tanto la clase fila de \(Z_n^{(r)}\) pertenece a
\(\operatorname{Gr}_{>0}(r,N_n+1)\). Las inserciones
\eqref{geometria_positiva:eq:carry-insertion} prolongan esta familia sin alterar el signo de
ningún menor antecedente.
\end{theorem}
\begin{proof}
El menor es un determinante de Vandermonde. Cada factor de
\eqref{geometria_positiva:eq:vandermonde-minor} es positivo por el orden estricto de los puntos.
Una inserción mantiene el orden de las columnas antiguas y añade puntos dentro
de sus intervalos. Los menores formados sólo por columnas antecedentes
conservan exactamente su valor; los que contienen una columna nueva vuelven a
ser productos de diferencias positivas.
\end{proof}

La versión ponderada conserva la misma orientación. Si
\(\lambda_{n,j}>0\), ponemos
\[
 C_{n;aj}=\lambda_{n,j}t_{n,j}^{a-1}.
\]
Entonces
\begin{equation}
 \Delta_I(C_n)=\left(\prod_{j\in I}\lambda_{n,j}\right)
 \Delta_I(Z_n^{(r)})>0.
 \label{geometria_positiva:eq:weighted-vandermonde}
\end{equation}
Los pesos pueden recibir una medida cilíndrica compatible o la raíz de una
multiplicidad de ruta, siempre que se declaren antes de evaluar el objeto
externo.

\subsection{Máscaras de supervivencia y realización positroidal}

El TPK conserva en cada profundidad una máscara de supervivencia
\(S_n\subseteq\{0,\ldots,N_n\}\). Para hacerla visible sin cambiar el
generador, se permite
\[
 \lambda_{n,j}>0\quad(j\in S_n),\qquad
 \lambda_{n,j}=0\quad(j\notin S_n).
\]

\paragraph{Estatuto.}
La asignación de una matriz de momentos a una máscara de supervivencia es una
\emph{formalización nueva}: utiliza datos ya generados por el TPK, pero no se
presenta como teorema previamente contenido en el corpus. Su contenido
matemático exacto es el siguiente.

\begin{proposition}[Realización positroidal de una máscara]
\label{geometria_positiva:thm:tpk-positroid-cell}
Supóngase \(|S_n|\ge r\). La matriz ponderada satisface
\[
 \Delta_I(C_{n,S_n})>0\quad\Longleftrightarrow\quad I\subseteq S_n,
 \qquad
 \Delta_I(C_{n,S_n})=0\quad\Longleftrightarrow\quad I\not\subseteq S_n.
\]
Su matroide es el matroide uniforme de rango \(r\) sobre \(S_n\), con las
posiciones complementarias como lazos; en particular, es un positroide. Las
realizaciones de dos niveles son compatibles siempre que sus máscaras lo sean
bajo el mapa de enlace y que las columnas antiguas conserven su identificación.
\end{proposition}
\begin{proof}
La factorización \eqref{geometria_positiva:eq:weighted-vandermonde} muestra que un menor se anula
exactamente cuando contiene un peso nulo. Sobre \(S_n\), todos los factores son
positivos y todos los subconjuntos de cardinal \(r\) son bases. Un matroide
uniforme realizado por puntos ordenados en la curva de momentos es positroidal;
añadir lazos conserva esa clase. La compatibilidad procede de que una máscara
refinada hereda los supervivientes y añade únicamente las columnas permitidas
por la hipótesis de compatibilidad.
\end{proof}

La proposición certifica la realización de cada máscara dentro del funtor de
momentos ordenados. No certifica todavía un atlas positroidal intrínseco del
TPK ni la cobertura de todas las celdas de
\(\operatorname{Gr}_{\ge0}(r,N_n+1)\). Para esa promoción se requiere una red
planar o medida de frontera derivada del TPK, junto con su permutación decorada
o collar de Grassmann, y la naturalidad de esos datos bajo refinamiento. La
ausencia de cualquiera de esas piezas falsaría la promoción global, no la
proposición anterior.

\subsection{Intercambios de Plücker y cartas cluster}

Para \(r=2\), escribamos
\(z_j=\lambda_j(1,t_j)^{\mathsf T}\). Si \(a<b<c<d\), se tiene
\begin{equation}
 \Delta_{ac}\Delta_{bd}
 =\Delta_{ab}\Delta_{cd}+\Delta_{ad}\Delta_{bc}.
 \label{geometria_positiva:eq:ptolemy-tpk}
\end{equation}

\paragraph{Estatuto.}
La relación siguiente es exacta para la realización de momentos ordenados. Su
lectura como estructura cluster del TPK es una \emph{formalización nueva}; no
identifica el acarreo con una mutación ni afirma que el corpus haya construido
ya un grafo plábico universal.

\begin{proposition}[Intercambio positivo en rango dos]
\label{geometria_positiva:thm:tpk-cluster-rank2}
Todas las coordenadas de \eqref{geometria_positiva:eq:ptolemy-tpk} son positivas y la mutación
que despeja cualquiera de las dos diagonales es sustracción libre. La inserción
de un hijo por \eqref{geometria_positiva:eq:carry-insertion} conmuta con la evaluación de la
relación de Ptolomeo. Por tanto, las triangulaciones del polígono ordenado y sus
flips proporcionan cartas cluster positivas para la realización ordenada de
cada nivel.
\end{proposition}
\begin{proof}
La identidad se obtiene sustituyendo
\(\Delta_{ij}=\lambda_i\lambda_j(t_j-t_i)\) y expandiendo. Todos sus términos
son positivos. La inserción conserva las coordenadas antiguas y asigna a la
nueva columna el parámetro ordenado exacto; evaluar primero el determinante o
primero la relación de intercambio produce la misma identidad polinómica.
\end{proof}

En rango general, para \(|I|=r-2\) y \(a<b<c<d\) fuera de \(I\),
\begin{equation}
 \Delta_{Iac}\Delta_{Ibd}
 =\Delta_{Iab}\Delta_{Icd}+\Delta_{Iad}\Delta_{Ibc}.
 \label{geometria_positiva:eq:short-plucker-tpk}
\end{equation}
Todos los términos son positivos en la celda superior y no negativos en las
celdas de supervivencia. Así, cada intercambio corto de Plücker es compatible
con el refinamiento dentro de esta realización. La afirmación correcta es que
el acarreo inserta datos y las mutaciones cambian de carta sobre esos datos; el
paso de acarreo no se identifica con un flip. Una compatibilidad cluster
intrínseca del TPK exige además derivar una semilla, su quiver o grafo plábico
y la acción de los mapas de enlace sobre las mutaciones.

\subsection{Formas canónicas y cancelación de fronteras interiores}

Para un intervalo orientado \(I=[a,b]\), su forma canónica es
\begin{equation}
 \Omega_I=\frac{dt}{t-a}-\frac{dt}{t-b}.
 \label{geometria_positiva:eq:interval-canonical-form}
\end{equation}
Si \(a=c_0<c_1<\cdots<c_m=b\), entonces
\begin{equation}
 \sum_{q=0}^{m-1}\Omega_{[c_q,c_{q+1}]}=\Omega_{[a,b]}.
 \label{geometria_positiva:eq:interval-telescopy}
\end{equation}
Cada polo interior aparece dos veces, con residuos opuestos.

La identidad unidimensional es exacta y no depende de una interpretación
positroidal adicional. Su extensión cubical se formula a continuación como una
\emph{formalización nueva exacta}. Sea
\(Q=\prod_{\ell=1}^d[a_\ell,b_\ell]\) con orientación
\(dx_1\wedge\cdots\wedge dx_d\). Definimos
\begin{equation}
 \Omega_Q=\bigwedge_{\ell=1}^d
 \left(\frac{dx_\ell}{x_\ell-a_\ell}
       -\frac{dx_\ell}{x_\ell-b_\ell}\right).
 \label{geometria_positiva:eq:box-canonical-form}
\end{equation}

\begin{theorem}[Formalización nueva: telescopía cubical multidimensional]
\label{geometria_positiva:thm:tpk-higher-residues}
Subdivídase cada intervalo coordenado por los puntos del refinamiento TPK y
denótense por \(Q_{\mathbf j}\) los cuboides resultantes, con la orientación
inducida. Entonces
\begin{equation}
 \boxed{\displaystyle
 \sum_{\mathbf j}\Omega_{Q_{\mathbf j}}=\Omega_Q.}
 \label{geometria_positiva:eq:box-telescopy}
\end{equation}
Todos los polos de facetas interiores se cancelan. Si
\(F_1\supset F_2\supset\cdots\supset F_s\) es una bandera de caras orientadas,
entonces
\begin{equation}
 \operatorname{Res}_{F_s}\cdots\operatorname{Res}_{F_1}\Omega_Q
 =\operatorname{sgn}(F_\bullet)\,\Omega_{F_s}.
 \label{geometria_positiva:eq:iterated-residue}
\end{equation}
En particular, el residuo en cada vértice tiene módulo uno y el signo fijado
por la orientación.
\end{theorem}
\begin{proof}
La suma sobre la cuadrícula es el producto exterior de las \(d\) sumas
unidimensionales. Aplicar \eqref{geometria_positiva:eq:interval-telescopy} en cada factor da
\eqref{geometria_positiva:eq:box-telescopy}. Una faceta interior tiene dos celdas incidentes y
orientaciones normales opuestas, de modo que sus residuos se cancelan. Tomar
un residuo elimina el factor normal correspondiente y deja la forma canónica
de la cara, con el signo de orientación. La iteración prueba
\eqref{geometria_positiva:eq:iterated-residue}.
\end{proof}

\begin{corollary}[Independencia del refinamiento]
La forma global no depende de cuántos niveles de acarreo se hayan insertado.
Refinar cambia la triangulación o cubiculación y preserva la suma de formas.
La memoria TPK conserva qué celdas intermedias se atravesaron, aunque la forma
global cancele sus fronteras internas.
\end{corollary}

\begin{proposition}[Criterio orientado general de cancelación]
\label{geometria_positiva:prop:oriented-residue-chain}
Sea \(X=\bigcup_\sigma X_\sigma\) una subdivisión positiva pura y orientada.
Supóngase que cada celda posee una forma \(\Omega_\sigma\) y que, para toda
faceta \(F\subset\overline{X_\sigma}\),
\[
\operatorname{Res}_F\Omega_\sigma=[\sigma:F]\Omega_F.
\]
Para una cadena \(c=\sum_\sigma w_\sigma[\sigma]\),
\begin{equation}
\operatorname{Res}_F\!\left(\sum_\sigma w_\sigma\Omega_\sigma\right)
=(\partial c)_F\Omega_F.
\label{geometria_positiva:eq:residue-chain-map}
\end{equation}
En particular, una faceta interior se cancela si y sólo si su coeficiente en
\(\partial c\) es cero. Los residuos iterados sobre banderas interiores se
cancelan por \(\partial^2=0\).
\end{proposition}
\begin{proof}
La linealidad del residuo da
\[
\operatorname{Res}_F\!\left(\sum_\sigma w_\sigma\Omega_\sigma\right)
=\sum_{\sigma\supset F}w_\sigma[\sigma:F]\Omega_F,
\]
y el coeficiente de \(\Omega_F\) es precisamente el de \(F\) en
\(\partial c\). La iteración identifica la suma de incidencias de codimensión
dos con el coeficiente correspondiente de \(\partial^2c\), que es cero.
\end{proof}

La proposición es una formalización nueva exacta. Para identificar la suma con
la forma canónica única de \(X\) debe excluirse además una forma superior
regular ambigua. El corpus ya aporta el complejo de caminos y memoria con
\(D^2=0\); la promoción positroidal requiere construir el morfismo
\[
\Omega_{\rm TPK}:C_\bullet^{\rm TPK}\longrightarrow
\Omega^\bullet_{\log},
\qquad
\operatorname{Res}\circ\Omega_{\rm TPK}
=\Omega_{\rm TPK}\circ D.
\]
En cilindros y productos este morfismo queda dado por
\eqref{geometria_positiva:eq:interval-canonical-form} y \eqref{geometria_positiva:eq:box-canonical-form}. Para un
atlas positroidal general deben añadirse estratos, incidencias, orientaciones,
unicidad y compatibilidad Plücker--cluster. Los pesos probabilísticos de las
historias tampoco pueden multiplicar indiscriminadamente las formas: una
faceta compartida conservaría un residuo proporcional a la diferencia de
pesos. Los factores \(\sqrt\mu\) de una matriz positiva pertenecen a otra
operación.

\subsection{El registro \(K\) como marco métrico y archivístico}

El registro dodecafásico se produce por
\[
 x_{\rm term}\mapsto\mathcal R_{12}
 \mapsto(b^{90},b^{120},Q_{\rm TPK})
 \mapsto U_{\rm sgn}\mapsto K
\]
y vale
\[
 K=(234,543,140,729,659,824,621,58,914,794,146,601).
\]
Sea \(S\) el desplazamiento cíclico de doce posiciones y
\[
 O_K=(K,SK,\ldots,S^{11}K),\qquad G_K=O_K^*O_K.
\]
El propietario demuestra
\begin{equation}
 \det O_K=5483119259214020183850397930008283625>0,
 \qquad
 589609I\le G_K\le6263^2I.
 \label{geometria_positiva:eq:K-frame-bounds}
\end{equation}
En particular, \(O_K\) es invertible y define un marco cuantitativamente
controlado. Como su determinante es positivo, su descomposición polar puede
escribirse
\[
 O_K=F_KP_K,\qquad P_K=G_K^{1/2}>0,\qquad F_K\in SO(12).
\]

\begin{proposition}[Covariancia exacta de los datos de Plücker]
\label{geometria_positiva:thm:K-plucker-transport}
Sea \(X\in\operatorname{Mat}_{r\times12}(\mathbb R)\) de rango \(r\) y
\(Y=XO_K^{\mathsf T}\). Entonces
\begin{equation}
 \Delta_I(Y)=
 \sum_{\substack{J\subseteq\{1,\ldots,12\}\\ |J|=r}}
 \Delta_J(X)\det\!\bigl((O_K)_{I,J}\bigr).
 \label{geometria_positiva:eq:K-cauchy-binet}
\end{equation}
Equivalentemente, el vector de Plücker se transforma por
\(\bigwedge^rO_K\). En rango doce,
\[
 \Delta_{\{1,\ldots,12\}}(XO_K^{\mathsf T})
 =(\det O_K)\Delta_{\{1,\ldots,12\}}(X).
\]
\end{proposition}
\begin{proof}
La fórmula \eqref{geometria_positiva:eq:K-cauchy-binet} es la identidad de Cauchy--Binet aplicada
a cada conjunto de columnas de \(Y\). La formulación exterior y el caso de
rango doce son sus expresiones functorial y maximal, respectivamente.
\end{proof}

La covariancia de Plücker no equivale a preservación de la positividad total
estándar. En efecto, \(O_K\) posee menores de ambos signos. Dos testigos
explícitos son
\begin{equation}
 \det\begin{pmatrix}234&543\\543&140\end{pmatrix}=-262089,
 \qquad
 \det\begin{pmatrix}234&140\\543&729\end{pmatrix}=94566.
 \label{geometria_positiva:eq:K-mixed-minors}
\end{equation}
Por ejemplo, la imagen por \(O_K\) del plano coordenado positivo generado por
los dos primeros vectores canónicos ya contiene el primer menor negativo de
\eqref{geometria_positiva:eq:K-mixed-minors}. Por tanto,
\begin{equation}
 \left\{[XO_K^{\mathsf T}]:
 [X]\in\operatorname{Gr}_{\ge0}(r,12)\right\}
 \not\subseteq\operatorname{Gr}_{\ge0}(r,12)
 \quad\text{en general}.
 \label{geometria_positiva:eq:K-not-standard-totally-positive}
\end{equation}
Así, \(O_K\) es un marco métrico y archivístico invertible; no es una matriz
de medida de frontera ni un generador de positividad total estándar.

\paragraph{Carta positiva adaptada a \(K\).}
La imagen exacta
\begin{equation}
 \operatorname{Gr}_{\ge0}^{K}(r,12)
 :=\left\{[XO_K^{\mathsf T}]:
 [X]\in\operatorname{Gr}_{\ge0}(r,12)\right\}
 \label{geometria_positiva:eq:K-adapted-positive-chart}
\end{equation}
define una carta positiva adaptada al marco. Si
\(\Phi_K([X])=[XO_K^{\mathsf T}]\), sus coordenadas adaptadas son
\begin{equation}
 \Delta_I^K([Y])
 :=\Delta_I\!\left(Y(O_K^{\mathsf T})^{-1}\right)\ge0.
 \label{geometria_positiva:eq:K-adapted-plucker}
\end{equation}
Esta positividad es exacta por transporte de estructura; no debe confundirse
con la positividad de los menores de \(Y\) en la base ambiente estándar. Del
mismo modo, una forma canónica \(\Omega\) en la carta original se transporta
como \((\Phi_K)_*\Omega\).

\paragraph{Formalización nueva: acción diagonal positiva de \(K\).}
Como todas las componentes \(K_i\) son estrictamente positivas, se define
\begin{equation}
 D_K(s):=\operatorname{diag}
 \bigl(e^{s\log K_1},\ldots,e^{s\log K_{12}}\bigr),
 \qquad s\in\mathbb R,
 \label{geometria_positiva:eq:K-positive-diagonal-flow}
\end{equation}
de modo que \(D_K(1)=\operatorname{diag}(K_1,\ldots,K_{12})\). Para toda
matriz \(X\in\operatorname{Mat}_{r\times12}\),
\begin{equation}
 \Delta_I(XD_K(s))
 =e^{\,s\sum_{i\in I}\log K_i}\Delta_I(X).
 \label{geometria_positiva:eq:K-diagonal-minors}
\end{equation}
Esta acción preserva exactamente
\(\operatorname{Gr}_{\ge0}(r,12)\) y su estratificación por anulaciones de
menores. La acción diagonal es una formalización nueva construida a partir del
registro generado \(K\); no se atribuye al propietario original de \(O_K\) ni
se utiliza para ocultar los signos mixtos de \eqref{geometria_positiva:eq:K-mixed-minors}.

\paragraph{Formalización nueva: flujo centrado determinado por \(K\).}
El corpus construye
\begin{equation}
 u_K=P_3P_{11}K,\qquad
 \|u_K\|^2=
 \frac{6638585+2275584\sqrt5}{20}>0,
 \qquad u_K\perp\mathbf1.
 \label{geometria_positiva:eq:K-centered-direction}
\end{equation}
Sea \(\widehat u_K=u_K/\|u_K\|\) y
\(\rho_A=\pi_{\rm HMT}/\varphi_{\rm HMT}^2>0\). Entonces
\begin{equation}
 g_{K,\rho}(\tau)=
 \operatorname{diag}\!\left(
 \rho_A^{\tau\widehat u_{K,1}},\ldots,
 \rho_A^{\tau\widehat u_{K,12}}\right)
 \in SL(12,\mathbb R)_{>0}.
 \label{geometria_positiva:eq:K-centered-positive-flow}
\end{equation}
Para toda matriz \(X\),
\begin{equation}
 \Delta_I(Xg_{K,\rho}(\tau))
 =\rho_A^{\tau\sum_{i\in I}\widehat u_{K,i}}\Delta_I(X).
 \label{geometria_positiva:eq:K-centered-flow-minors}
\end{equation}
Por tanto el flujo preserva signos y menores nulos. El registro \(K\) fija la
dirección y la razón generada \(\rho_A\) normaliza la velocidad. La
razón \(\rho_A\) se evalúa aquí en una etapa posterior: \(\pi_{\rm HMT}\) y
\(\varphi_{\rm HMT}\) son dos coordenadas de la cuatrirrelación
\((\pi,\varphi,e,\alpha)\) generada coinductivamente por el mismo estado
APP--TRIT--TPK\@. Ni sus valores convencionales ni la razón ya evaluada
seleccionan \(K\), \(u_K\) o el refinamiento. La
identificación de \(\tau\) con tiempo físico exige un reloj \(\tau(t)\); la
prolongación a toda profundidad exige un cociclo sectorial compatible con
\(j\mapsto Bj+c\).

El archivo completo
\[
 \mathcal Vv=\left(A_\infty^{1/2}v,(D_jR_jv)_{j\ge0}\right),
 \qquad\mathcal V^*\mathcal V=I,
\]
conserva esta información en todos los grados:
\begin{equation}
 \left\langle\bigwedge^r\mathcal V\,\xi,
 \bigwedge^r\mathcal V\,\eta\right\rangle
 =\langle\xi,\eta\rangle,
 \qquad
 \|\wedge^r\mathcal V\,\xi\|=\|\xi\|.
 \label{geometria_positiva:eq:exterior-isometry}
\end{equation}
Por tanto conserva determinantes de Gram y normas de Plücker. Esta isometría
no afirma que \(\mathcal V\) convierta los menores de \(O_K\) en menores
positivos: preserva la información archivada en el espacio y la orientación
que se transporten explícitamente.

La lectura racional
\[
 \kappa_{\rm ag}=
 \frac{234543140729659824621058914794146601}{10^{36}-1}
\]
es reversible una vez fijadas base, longitud y origen de fase. La función de
acoplamiento de \(K\) se expresa, por ello, en tres capas distintas: su lectura
aritmética \(\kappa_{\rm ag}\), su Gram positivo \(G_K\) y su transporte
incidencial. Las tres son publicaciones del registro generado; ninguna se usa
para elegir retrospectivamente la historia.

\subsection{Flujo areal--volumétrico y razón marcada}

El calendario TPK induce la matriz primitiva
\[
 F_{\rm av}=\begin{pmatrix}0&1\\1&1\end{pmatrix},
\]
cuyas potencias obedecen la recurrencia de Fibonacci. La medida de Parry de su
envolvente simbólico distingue los modos areal y volumétrico y satisface
\begin{equation}
 \frac{\mu_{\rm ar}}{\mu_{\rm vol}}=\varphi_{\rm HMT}^{-2}.
 \label{geometria_positiva:eq:area-volume-ratio}
\end{equation}
La lectura regional marca una sola vez esa razón:
\begin{equation}
 \Theta_n=\pi_{\rm HMT}\frac{F_{n-1}}{F_{n+1}}
 \longrightarrow\frac{\pi_{\rm HMT}}{\varphi_{\rm HMT}^2}.
 \label{geometria_positiva:eq:marked-area-volume}
\end{equation}
El lector recíproco
\(R(x,y)=x/y^2\), conjugado por \(J(x,y)=(y,x)\), produce la orientación
\(\varphi_{\rm HMT}/\pi_{\rm HMT}^2\). Estas razones son lectores tipados del
flujo; no se identifican con \(\kappa_{\rm ag}\), aunque todas participan en
la misma arquitectura de acoplamiento.

La forma canónica del complejo es estable bajo el refinamiento, mientras
\eqref{geometria_positiva:eq:area-volume-ratio} determina la proporción estacionaria de sus dos
modos. Éste es el cierre del círculo: el refinamiento aritmético genera la
subdivisión; la razón holográfica gobierna su distribución asintótica; \(K\)
conserva la métrica y la orientación; y la forma canónica retiene sólo la
frontera externa después de cancelar las paredes internas.

\subsection{Centro electrónico y sección material}

Normalicemos el atlas \(9\times9\) por
\[
 \nu(r,s)=\left(\frac{r-1}{8},\frac{s-1}{8}\right).
\]
La inversión celular \(\rho_c(r,s)=(10-r,10-s)\) se convierte en
\((u,v)\mapsto(1-u,1-v)\). Su único punto fijo es
\begin{equation}
 \nu(5,5)=\left(\frac12,\frac12\right).
 \label{geometria_positiva:eq:electron-center}
\end{equation}
Los \(81\) estados constituyen puntos de una malla cuya cubiculación elemental
posee \(64\) rectángulos; no se presentan como \(81\) celdas bidimensionales.

\begin{theorem}[Unicidad de la sección puntual equivariante]
\label{geometria_positiva:thm:electron-equivariant-section}
En una fibra estable sobre la que la etiqueta gruesa permanece fija, toda
selección puntual equivariante para \(\rho_c\) debe tomar el valor \((5,5)\).
La fibra electrónica de nivel cero contiene ese único punto; ninguna otra
fibra admite tal selección.
\end{theorem}
\begin{proof}
Si \(a(y)\) es equivariante y \(y\) es fijo, entonces
\(a(y)=a(\rho_cy)=\rho_ca(y)\). La única solución de
\(\rho_c(r,s)=(r,s)\) es \(r=s=5\). El censo del atlas sitúa esa celda en la
fibra electrónica de nivel cero.
\end{proof}

El teorema explica por qué el electrón es una sección puntual central y las
demás especies se presentan mediante órbitas o multisecciones. En particular,
deben mantenerse separados los siguientes objetos tipados:
\begin{enumerate}
 \item la celda fija \((5,5)\), perteneciente al atlas \(9\times9\);
 \item el dato \(111111\), perteneciente a la quinta frontera nonádica;
 \item la firma multiplicativa transversal
       \(\mathcal E_{\times}=555555\), perteneciente al lector regional de
       \(\pi\);
 \item el bloque de signo \((+1,+1,+1)\), que aparece dentro de la palabra
       dodecafásica
       \[
        \tau_e=(+1,+1,+1,-1,-1,-1,+1,+1,+1,-1,-1,-1).
       \]
\end{enumerate}
La estructura APP central, distinta de esos tres registros, es
\[
 \mathcal A(5,5)=((1,1),(7,2)),\qquad
 5+5=1+9\cdot1,\qquad 5\cdot5=7+9\cdot2.
\]
La coincidencia o repetición de dígitos no constituye una identificación:
\((5,5)\), \(111111\) y \(555555\) poseen dominios, operaciones y funciones
causales diferentes.

La positividad material es posterior. Si \(M_M^{(0)}\ge0\) y
\(R_{\rm act}>0\), entonces
\begin{equation}
 M_M^{\beta,r}=R_{\rm act}^{B_M^r/2}
 M_M^{(0)}R_{\rm act}^{B_M^r/2}\ge0.
 \label{geometria_positiva:eq:material-positive-congruence}
\end{equation}
La congruencia conserva el cono positivo y realiza sobre la rama de masa la
positividad previamente estructurada; no convierte la masa observada en una
entrada del atlas.

\subsection{Multimorfología positiva del estado enriquecido}

La multimorfología no designa una superposición informal de figuras. Es una
familia de morfismos tipados con fuente común:
\[
\begin{aligned}
\mathfrak P_n&:\widehat{\mathcal X}_{\rm TPK,n}
\longrightarrow\operatorname{Gr}_{\ge0}(k,N_n),\\
\mathfrak R_n&:\widehat{\mathcal X}_{\rm TPK,n}
\longrightarrow C^+_{d,n},\\
\mathfrak C_n&:\widehat{\mathcal X}_{\rm TPK,n}
\longrightarrow\mathcal Q^+_{\rm con,n},\\
\mathfrak M_n&:\widehat{\mathcal X}_{\rm TPK,n}
\longrightarrow\mathcal L^+_{M,n}.
\end{aligned}
\]
Sus imágenes son, respectivamente, las positividades de Plücker, reticular,
constitutiva y operatoria de masa. Cada morfismo conserva las coordenadas que
su lector necesita y conmuta con el enlace de refinamiento de su codominio.
Los cuatro codominios no se identifican ni se utilizan para seleccionarse
entre sí.

\begin{proposition}[Coherencia multimorfológica]
\label{geometria_positiva:prop:multimorphology-coherence}
Si cada morfismo conmuta con truncamiento, frontera y residuo, la cancelación
de caras interiores se transporta a las cuatro realizaciones. Si además una
aplicación cinemática
\[
\mathfrak A_{Z,n}:C\longmapsto CZ^{\mathsf T}
\]
es de grado uno, preserva orientación y cubre su imagen con fronteras
compatibles, la forma amplituhedral inducida es independiente de la
profundidad de representación.
\end{proposition}
\begin{proof}
La Proposición~\ref{geometria_positiva:prop:oriented-residue-chain} se aplica antes de cada
realización. Los cuadrados de compatibilidad transportan el coeficiente de
incidencia y el residuo. La condición de grado uno evita multiplicidades
interiores en el empuje cinemático.
\end{proof}

La realización \(\mathfrak P_n\) por curva de momentos es exacta como
formalización ordenada; su promoción a atlas positroidal intrínseco conserva
las condiciones ya declaradas. La rama reticular y la operatoria poseen
propietarios HMT--MD previos. La rama constitutiva debe especificar en cada
aplicación su forma positiva y su lector. La multimorfología reúne estos mapas;
no añade una quinta positividad.

% FORMALIZACION_REUNIDA: composición de la respuesta constitutiva de III
% con el lector positivo de GEOMETRIA_POSITIVA/08_ARTICULO_AUTONOMO.tex.
% Propietario de las operaciones: III/manuscrito__sections__04_respuesta_constitutiva.tex.
% Residencia propuesta: después de la coherencia multimorfológica, antes del
% empuje amplituhedral. No modifica sus hipótesis de frontera, residuo o grado.
\subsection{Realización constitutiva positiva y naturalidad espectral}
\label{empalme:positividad-constitutiva}

La realización constitutiva de la
proposición~\ref{geometria_positiva:prop:multimorphology-coherence} recibe
un operador ya construido desde las dos hojas orientadas del TPK. La
sección~\ref{iii:iii:sec:constitutiva} determina su respuesta y su inversa;
la presente composición explicita la forma positiva y el enlace de
refinamiento que utiliza esa realización.

Sea $\widehat{\mathcal X}_{\rm con}$ el dominio de estados enriquecidos
cuya lectura angular pertenece a la cámara $x>y>0$. Para
$\xi\in\widehat{\mathcal X}_{\rm con}$, la lectura produce un espacio
de dos hojas $E_\xi$, una involución ortogonal $\mathcal R_\xi$ y el
transporte
\[
 P_{\pm,\xi}=\frac{I\pm\mathcal R_\xi}{2},\qquad
 T_\xi=q_{+,\xi}P_{+,\xi}+q_{-,\xi}P_{-,\xi},
 \qquad 0<q_{+,\xi}<q_{-,\xi}<1.
\]
Las incidencias $90$ y $120$ proceden del calendario y actúan sobre
ese mismo transporte. La respuesta es
\[
 V_\xi=(I-T_\xi^{90})(I-T_\xi^{120})^{-1},\qquad
 W_{m,\xi}=V_\xi\otimes I_{9^m}.
\]
La amplificación conserva todas las posiciones de la nueva fibra
nonádica. Sobre $E_{m,\xi}=E_\xi\otimes\mathbb C^{9^m}$ se define
la forma
\begin{equation}
 \mathfrak C_m(\xi)(v,w)
 :=\langle W_{m,\xi}v,W_{m,\xi}w\rangle
 =\langle v,W_{m,\xi}^{2}w\rangle.
 \label{empalme:eq:forma-constitutiva}
\end{equation}
Ésta es una realización concreta en el cono de formas hermíticas
positivas definidas. El cuadrado conserva las respuestas cuadráticas
de las dos hojas que publican las coordenadas
$\widehat\varepsilon$ y $\widehat\mu$ de
\eqref{iii:iii:eq:four}.

\begin{proposition}[Coercividad, recuperación y enlace nonádico]
\label{empalme:prop:constitutiva-natural}
La forma~\eqref{empalme:eq:forma-constitutiva} satisface, para todo
$v\in E_{m,\xi}$ no nulo,
\[
 \frac9{16}\|v\|^2
 <\mathfrak C_m(\xi)(v,v)<\|v\|^2.
\]
Si $A_{m,\xi}=W_{m,\xi}^{2}$ y $\mathbb E_m$ es la expectativa
parcial normalizada sobre el último factor de dimensión nueve, entonces
\[
 \mathbb E_m(A_{m+1,\xi})=A_{m,\xi}.
\]
La forma, junto con las etiquetas orientadas de hoja, recupera
$V_\xi$, los dos canales $q_{\pm,\xi}$ y sus coordenadas angulares.
Si un transporte unitario $J:E_\xi\to E_{\xi'}$ satisface
$JT_\xi=T_{\xi'}J$, entonces
\[
 (J\otimes I)^*A_{m,\xi'}(J\otimes I)=A_{m,\xi}.
\]
La identidad se mantiene para cualquier función continua del espectro
que se aplique a la respuesta antes de amplificar.
\end{proposition}

\begin{proof}
Los proyectores de hoja son ortogonales y suman la identidad. El cálculo
espectral de~\eqref{iii:iii:eq:rchannels} da
\[
 V_\xi=r_{+,\xi}P_{+,\xi}+r_{-,\xi}P_{-,\xi},\qquad
 r_{\pm,\xi}=f(q_{\pm,\xi}^{30}),\qquad
 f(s)=\frac{1+s+s^2}{1+s+s^2+s^3}.
\]
En $0<s<1$ se tiene
\[
 4(1+s+s^2)-3(1+s+s^2+s^3)
 =(1-s)(3s^2+2s+1)>0,
\]
y el denominador supera al numerador en $s^3>0$. Por tanto
$3/4<r_{\pm,\xi}<1$. La descomposición ortogonal y el cuadrado de
estos límites prueban la coercividad, que permanece idéntica bajo
amplificación.

Como $A_{m+1,\xi}=A_{m,\xi}\otimes I_9$ y la traza normalizada
del último factor satisface $\operatorname{tr}_9(I_9)=1$,
\[
 \mathbb E_m(A_{m+1,\xi})
 =(\operatorname{id}\otimes\operatorname{tr}_9)
 (A_{m,\xi}\otimes I_9)=A_{m,\xi}.
\]
La misma igualdad vale para cada polinomio en $W_{m,\xi}$. La
aproximación uniforme sobre el espectro compacto y la continuidad de
la expectativa la extienden a toda función continua.

La raíz positiva de $A_{0,\xi}$ recupera $V_\xi$. Los proyectores
orientados identifican sus dos autovalores. La función $f$ es
estrictamente decreciente en $(0,1)$ por
\eqref{iii:iii:eq:monotone}; su inversa determina las raíces únicas
de~\eqref{iii:iii:eq:cubic}. La raíz positiva de orden treinta
restituye $q_{\pm,\xi}$, y
\eqref{iii:iii:eq:recover-angular} restituye $x$ e $y$.

Finalmente, $JT_\xi=T_{\xi'}J$ implica la misma identidad para
todas sus potencias. La invertibilidad de $I-T^{120}$ proporciona
\[
 J(I-T_\xi^{120})^{-1}
 =(I-T_{\xi'}^{120})^{-1}J.
\]
De aquí resulta $JV_\xi=V_{\xi'}J$. La amplificación, el cuadrado
y la unitariedad de $J$ dan la congruencia declarada. El argumento
polinómico y su extensión uniforme prueban la última afirmación.
\end{proof}

El enlace de esta construcción es la expectativa normalizada que
conserva el operador amplificado. La compatibilidad con un transporte
geométrico utiliza el entrelazamiento $JT_\xi=T_{\xi'}J$ expresado
en el enunciado. Por su parte, el transporte de formas canónicas y
residuos mantiene los cuadrados de frontera de la
proposición~\ref{geometria_positiva:prop:multimorphology-coherence}.
Así quedan precisadas las operaciones de cada nivel: respuesta angular,
forma constitutiva positiva, amplificación de la fibra y transporte de
incidencias. La restitución angular recupera esta lectura; el archivo
enriquecido conserva además la historia que la produjo.

\subsection{Formalización amplituhedral condicionada}

\paragraph{Estatuto.}
Lo que sigue es una \emph{formalización nueva y condicional}. El corpus
proporciona el refinamiento, las rutas, la supervivencia y el dato ordenado;
la identificación de su imagen con un amplituhedro y la promoción de la suma
celular a forma canónica requieren las hipótesis geométricas que se enuncian
explícitamente.

Fijemos \(k,m\ge1\) y \(N=N_n+1\ge k+m\). Sea
\(Z_n=Z_n^{(k+m)}\) el dato externo positivo construido por el acarreo y sea
\(C_n\in\operatorname{Gr}_{\ge0}(k,N)\) una matriz de rutas supervivientes.
En el locus donde \(\operatorname{rank}(C_nZ_n^{\mathsf T})=k\), definimos
\begin{equation}
 Y_n=C_nZ_n^{\mathsf T}\in\operatorname{Gr}(k,k+m).
 \label{geometria_positiva:eq:hmt-amplituhedron-map}
\end{equation}

\begin{definition}[Imagen amplituhedral candidata]
Sea \(\mathfrak C_n^{\rm TPK}\) el complejo positivo de rutas y acarreos a
profundidad \(n\). Su imagen
\[
 \mathcal A_n^{\rm HMT}(k,m)
 =\{C Z_n^{\mathsf T}:C\in\mathfrak C_n^{\rm TPK}\}
\]
se denomina imagen amplituhedral candidata de nivel \(n\). Denotamos por
\[
 \mathcal F_{{\rm kin},n}:
 \mathfrak C_n^{\rm TPK}\longrightarrow
 \operatorname{Gr}(k,k+m),\qquad
 [C]\longmapsto[CZ_n^{\mathsf T}]
\]
su aplicación cinemática.
\end{definition}

\begin{proposition}[Empuje canónico bajo hipótesis de grado y frontera]
\label{geometria_positiva:thm:hmt-amplituhedron-form}
Supóngase que \(\mathfrak C_n^{\rm TPK}\) es un complejo orientado de dimensión
pura, que \(\mathcal F_{{\rm kin},n}\) es propia sobre su clausura,
genéricamente finita de grado uno en cada componente, preserva la orientación,
cubre su imagen sin solapamientos interiores no declarados y lleva fronteras a
fronteras de modo compatible con los residuos. Entonces
\begin{equation}
 \Omega_n^{\rm HMT}
 =(\mathcal F_{{\rm kin},n})_*
 \left(\sum_{\sigma\in\mathfrak C_n^{\rm TPK}}\Omega_\sigma\right)
 \label{geometria_positiva:eq:hmt-canonical-pushforward}
\end{equation}
está bien definida, tiene polos logarítmicos únicamente en la frontera externa
de \(\mathcal A_n^{\rm HMT}\), y es independiente de toda subdivisión ulterior
compatible. Sus residuos son las formas canónicas de las imágenes de las caras.
\end{proposition}
\begin{proof}
El teorema~\ref{geometria_positiva:thm:tpk-higher-residues} cancela las facetas interiores antes
del empuje. La propiedad de grado uno y la ausencia de solapamientos interiores
impiden multiplicidades espurias; la compatibilidad frontera--residuo transporta
los polos logarítmicos y sus residuos. Un refinamiento compatible cambia la
descomposición celular, pero no su suma.
\end{proof}

La proposición determina una forma estructural y recuperable únicamente cuando
se verifican sus hipótesis. Si se demuestra adicionalmente que
\(\mathfrak C_n^{\rm TPK}=\operatorname{Gr}_{\ge0}(k,N)\) en el dominio
relevante y que el mapa tiene la cobertura requerida, entonces
\(\mathcal A_n^{\rm HMT}(k,m)\) coincide con el amplituhedro estándar de esos
datos y \eqref{geometria_positiva:eq:hmt-canonical-pushforward} es su forma canónica completa. Sin
la derivación del complejo, la prueba de las hipótesis de empuje y la cobertura,
el objeto permanece como formalización candidata; no es un teorema atribuido
al corpus ni una afirmación de suryectividad.

\subsection{Síntesis}

La multimorfología estructural obtenida puede escribirse como la composición
\[
 \boxed{
 \begin{gathered}\text{historia TPK}\longmapsto\text{celda ordenada}\\
 \longmapsto\text{realización positroidal ordenada}\\
 \longmapsto\text{forma canónica}\longmapsto\text{archivo orientado por }K.\end{gathered}}
\]
Cada flecha conserva una información distinta: el TPK conserva la genealogía;
el Grassmanniano conserva los menores; la forma canónica conserva la frontera
después de cancelar paredes internas; \(K\) conserva el marco métrico y la
recuperación. La estructura discreta del continuo aparece así como un sistema
que admite simultáneamente refinamiento infinito, positividad finita en cada
nivel y un límite archivado sin pérdida de la historia que lo generó. La
realización positroidal, las cartas cluster, la extensión cubical y la imagen
amplituhedral son formalizaciones añadidas con los estatutos precisos
declarados en cada apartado; no se retroproyectan como teoremas ya demostrados
por el corpus.

% END OWNER


\section{Estabilidad por refinamiento y compatibilidad de los lectores sectoriales}

\begin{theorem}[Unidad y distribución de la aportación]
\label{geometria_positiva:thm:unidad-distribucion}
El teorema de estabilidad por refinamiento tiene residencia primaria en la
construcción conjunta del continuo. Sus aplicaciones a la forma de Weil, las
cinco realizaciones terminales, el acoplamiento aritmético--geométrico y el
operador de masa son corolarios tipados. Ninguna de esas aplicaciones puede
seleccionar retrospectivamente el refinamiento ni sustituir su prueba.
\end{theorem}
\begin{proof}
La forma canónica y su cancelación se definen antes de aplicar cualquier
lector sectorial. Cada lector recibe la misma suma celular y debe conmutar con
frontera y residuo. La propiedad común se demuestra una sola vez; la
particularidad de cada sector reside en su mapa de realización. Si un lector
seleccionase las celdas a partir de su resultado, invertiría la cadena causal.
\end{proof}

\section{Consecuencias estructurales}

\subsection{Multimorfología del Grassmanniano y del amplituhedro}

La palabra \emph{multimorfología} designa aquí una estructura matemática
precisa. Sea \(\widehat{\mathcal X}_{\rm TPK}\) el espacio de estados
enriquecidos y sea \(p_{n+1,n}\) su sistema de truncamientos. Un dato
multimorfológico positivo es una familia de realizaciones tipadas
\[
\begin{aligned}
\mathfrak P_n&:\widehat{\mathcal X}_{\rm TPK,n}
 \longrightarrow \operatorname{Gr}_{\geq0}(k,N_n),\\
\mathfrak R_n&:\widehat{\mathcal X}_{\rm TPK,n}
 \longrightarrow C^{+}_{d,n},\\
\mathfrak C_n&:\widehat{\mathcal X}_{\rm TPK,n}
 \longrightarrow \mathcal Q^{+}_{\rm con,n},\\
\mathfrak M_n&:\widehat{\mathcal X}_{\rm TPK,n}
 \longrightarrow \mathcal L^{+}_{M,n},
\end{aligned}
\]
donde los codominios representan, respectivamente, la positividad de
Plücker, la reticular, la constitutiva y la operatoria de masa. Los cuatro
mapas deben conmutar con sus enlaces de refinamiento y conservar las
coordenadas del estado que cada lectura utiliza. No se exige que los cuatro
codominios sean iguales ni que uno seleccione a los demás.

El mapa cinemático
\[
\mathfrak A_{Z,n}:C\longmapsto CZ^{\mathsf T}
\]
compone con \(\mathfrak P_n\) y produce la morfología amplituhedral. Así, el
Grassmanniano organiza las cartas, los estratos y sus mutaciones; el
amplituhedro organiza las imágenes cinemáticas y sus formas canónicas. La
multimorfología es el diagrama que mantiene simultáneamente esas realizaciones
como imágenes del mismo estado enriquecido.

\begin{proposition}[Coherencia multimorfológica por refinamiento]
\label{geometria_positiva:prop:coherencia-multimorfologica}
Si cada realización conmuta con el truncamiento y su mapa de formas conmuta
con frontera y residuo, la cancelación de caras interiores se transporta a
los cuatro codominios y a toda imagen amplituhedral de grado uno. Las formas
publicadas son independientes de la profundidad aunque sus archivos conserven
las rutas refinadas.
\end{proposition}
\begin{proof}
La forma del estado fuente es la suma orientada de sus celdas. La conmutación
con el truncamiento identifica las dos representaciones de cada refinamiento;
la conmutación con frontera y residuo transporta la cancelación de las
incidencias opuestas. Un mapa cinemático de grado uno preserva la multiplicidad
de la forma en cada interior. Por tanto sólo sobreviven las fronteras externas
en cada realización.
\end{proof}

La existencia de \(\mathfrak R_n\) y \(\mathfrak M_n\) se apoya en los
propietarios reticular y material ya construidos. La realización
\(\mathfrak P_n\) por curva de momentos es exacta como formalización ordenada;
su promoción a atlas positroidal intrínseco exige la medición de frontera. La
realización amplituhedral exige además rango, grado y cobertura. Ésta es la
frontera probatoria de la multimorfología, no una restricción de su definición.

\subsection{Continuo}

El paso al límite no añade puntos desconectados: estabiliza formas compatibles
de una misma estructura refinada. Las paredes internas desaparecen de la forma
global, pero permanecen en la memoria de historias. Se obtiene así una
distinción exacta entre invariancia geométrica y conservación genealógica.

\subsection{Riemann--Weil}

La telescopía hace independiente de la malla nonádica la realización de la
forma aritmética. La positividad de Weil continúa dependiendo de los canales
gamma, primo y polar y de la identidad de momentos mixtos que los convierte en
Gram. El resultado nuevo no reemplaza esa prueba: demuestra que ningún polo
interior del refinamiento puede introducir una dirección negativa espuria.

\subsection{Cinco realizaciones terminales}

Un lector terminal transporta la cancelación si conmuta con la frontera y si
el residuo de la forma transportada coincide con la forma de la cara. Esta
condición unifica el mecanismo de pegado sin disgregar el continuo ni convertir
los cinco problemas en cinco generadores.

\subsection{Acoplamiento y flujo}

La constante holográfica de acoplamiento aritmético--geométrico recupera
\(K\), y \(K\) determina \(O_K\), \(G_K\) y el archivo. La acción general
\(D_K(s)\) abre una vía positiva. La dirección centrada
\(u_K=P_3P_{11}K\), que satisface \(u_K\perp\mathbf1\), permite afinarla a
\[
g_{K,\rho}(\tau)=
\operatorname{diag}
\bigl(\rho_A^{\tau\widehat u_{K,1}},\ldots,
\rho_A^{\tau\widehat u_{K,12}}\bigr)\in SL(12,\mathbb R)_{>0}.
\]
Las razones
\(\varphi^{-2}\), \(\pi/\varphi^2\) y \(\varphi/\pi^2\) pueden parametrizar
secciones diferentes del flujo después de ser generadas y tipadas como
publicaciones de la cuatrirrelación coinductiva
\((\pi,\varphi,e,\alpha)\); ninguna selecciona retrospectivamente el TPK ni se
confunde con \(\kappa_{\rm ag}\).

\subsection{Partículas y masas}

La congruencia
\[
M_M^{\beta,r}
=R_{\rm act}^{B_M^r/2}M_M^{(0)}R_{\rm act}^{B_M^r/2}
\]
preserva el cono positivo. La excepción electrónica es anterior a esta
evaluación: procede del único punto fijo \((5,5)\). La frontera \(111111\) y
la firma \(555555\) permanecen como registros distintos de la misma genealogía.

\section{Criterio de promoción a atlas positroidal y amplituhedro}

La promoción queda cerrada cuando se construyan materialmente:
\begin{enumerate}
\item una red planar \(\mathcal N_z\) por cilindro, con pesos positivos
derivados del TPK;
\item la matriz de medida de frontera y su fórmula LGV para todos los menores;
\item la permutación decorada o collar de Grassmann de cada estrato;
\item una semilla o grafo plábico cuya mutación conmute con el refinamiento y
la memoria;
\item el mapa cinemático \(C\mapsto CZ^{\mathsf T}\), con rango, dimensión,
grado, cobertura y correspondencia de fronteras;
\item la identidad de empuje de formas canónicas.
\end{enumerate}
Estas condiciones no son una lista sociológica: son los falsadores que
distinguen una realización completa de una colección de puntos positivos.

\section{Conclusión}

La aportación central es una ley de estabilidad geométrica del refinamiento:
la profundidad añade memoria y resolución, mientras la forma canónica conserva
la frontera total. El registro \(K\) hace recuperable y métrico ese proceso;
su dirección centrada organiza un flujo de determinante uno que respeta los
signos de Plücker. El electrón ocupa el único centro equivariante de la rama material.
Con estas piezas, el paso al Grassmanniano positivo deja de ser una analogía:
posee una realización ordenada exacta y un programa de promoción cuyos mapas y
falsadores están completamente especificados.

